\section{SHA-256 Cryptographic Analysis}
\label{sec:sha256}

\subsection{Algorithm Specification}

SHA-256 is a member of the SHA-2 family of cryptographic hash functions,
standardised by NIST in FIPS PUB 180-4~\citep{fips180}.
It maps an arbitrary-length input to a 256-bit (32-byte) output.

Formally, SHA-256 satisfies three properties~\citep{bellare2004hash}:
\begin{enumerate}
\item \textbf{Pre-image resistance}: Given a hash output $h$, it is
  computationally infeasible to find an input $m$ such that
  $\mathrm{SHA256}(m) = h$.
\item \textbf{Second pre-image resistance}: Given an input $m_1$, it is
  computationally infeasible to find a different input $m_2 \neq m_1$ such
  that $\mathrm{SHA256}(m_1) = \mathrm{SHA256}(m_2)$.
\item \textbf{Collision resistance}: It is computationally infeasible to
  find any two distinct inputs $m_1 \neq m_2$ such that
  $\mathrm{SHA256}(m_1) = \mathrm{SHA256}(m_2)$.
\end{enumerate}

\subsection{The Birthday Bound}

By the birthday paradox, a collision attack on a hash function with
$n$-bit output requires approximately $2^{n/2}$ hash evaluations in
expectation.
For SHA-256, this is $2^{128}$ evaluations.

However, the relevant security property for the audit chain is second
pre-image resistance (Property 2), not collision resistance.
The adversary has a specific target: the legitimate file with its known
hash.
They need to find a \emph{different} file that hashes to the same value.
Second pre-image attacks are harder than collision attacks; the best known
bound is $2^{256}$ for pre-image and $2^{128}$ for second pre-image
under the generic attack model.

For audit chain purposes (does this file match this hash?), we use
the collision resistance bound as a conservative estimate:
$2^{128}$ operations.
Under the birthday bound interpretation (finding any collision among
$2^{80}$ attempted values with 50\% probability), the bound is $2^{80}$.

\begin{theorem}[SHA-256 Audit Chain Security]
Under SHA-256, an adversary cannot produce a modified input file that
has the same SHA-256 hash as the original without performing at least
$2^{80}$ hash evaluations (birthday bound for collision finding).
At $10^{15}$ evaluations per second on 2026-grade hardware, this requires
approximately $3 \times 10^7$ years.
\end{theorem}

\subsection{Known Attack Status}

As of 2026, no practical attacks on SHA-256 are known.
The best theoretical attacks (differential cryptanalysis, reduced-round
attacks) apply only to reduced versions of SHA-256 and do not reduce
the security below $2^{128}$~\citep{mendel2013rebound}.
SHA-256 is approved for use in U.S.\ government applications under
NIST SP 800-57 through at least 2030.

\subsection{Quantum Resistance}

Grover's algorithm~\citep{grover1996fast} provides a quadratic speedup
for pre-image attacks on hash functions using a quantum computer,
reducing the effective security from $2^{256}$ to $2^{128}$.
For the collision resistance bound relevant to the audit chain, Grover's
algorithm provides no improvement (collision-finding does not reduce
below $2^{128}$ even with quantum speedup).

The \bisect{} audit chain is therefore quantum-resistant at the
$2^{128}$ security level for the collision-finding task.
